Evaluating large language models (LLMs) remains a critical bottleneck: benchmarks saturate rapidly (MMLU accuracy rose from 43.9\% to over 90\% in four years), human evaluation is expensive and non-reproducible, and existing automated metrics lack formal mathematical guarantees regarding stability, composability, or drift detection.
We present a comprehensive survey of LLM evaluation across 11 dimensions, identify 6 cross-cutting gaps in current methodology, and propose 5 novel mathematically grounded evaluation frameworks that bridge the empirical AI evaluation community with the mathematical sciences.
Our frameworks draw on persistent homology, sheaf theory, information geometry, topological data analysis, and random matrix theory to address fundamental limitations in existing approaches.
We validate each framework through proof-of-concept implementations on MMLU data (14,042 questions, 57 subjects, 12 model profiles).
Key results include: (1) topological drift detection with formal stability bounds (Wasserstein distance 1.77 between model generations), (2) sheaf-theoretic composition revealing when evaluation aggregation loses information (consistency ratios ranging from 0.48 to 1.83), (3) Fisher-Rao geometry providing proper metric structure for model comparison (maximum distance 2.82), (4) failure mode landscape analysis discovering phase transitions in error structure, and (5) spectral contamination detection via Marchenko-Pastur bounds (Tracy-Widom statistic 83.1).
These frameworks provide the first formal mathematical guarantees for LLM evaluation, opening new directions for rigorous, stable, and composable assessment of language models.
